% swim-times-srs.tex
% Software Requirements Specification for the swim-times program
% IEEE/ISO/IEC/IEEE 29148-2018 structure
% Compiled with: pdflatex swim-times-srs.tex

\documentclass[12pt]{article}

\usepackage{geometry}
\geometry{letterpaper, margin=1in}

\usepackage{hyperref}
\usepackage{booktabs}
\usepackage{array}
\usepackage{longtable}
\usepackage{enumitem}
\usepackage{titlesec}
\usepackage{fancyhdr}
\usepackage{xcolor}

\definecolor{castblue}{RGB}{0,48,135}

\pagestyle{fancy}
\fancyhf{}
\renewcommand{\headrulewidth}{0.4pt}
\fancyhead[L]{\small College Area Swim Team}
\fancyhead[R]{\small SRS: swim-times}
\fancyfoot[C]{\thepage}

\titleformat{\section}{\large\bfseries\color{castblue}}{\thesection}{0.5em}{}[\titlerule]
\titleformat{\subsection}{\normalsize\bfseries}{\thesubsection}{0.5em}{}

\newcommand{\code}[1]{\texttt{#1}}
\newcommand{\req}[1]{\textbf{[#1]}}

\begin{document}

% -----------------------------------------------------------------
% Title block
% -----------------------------------------------------------------
\begin{titlepage}
  \centering
  \vspace*{2cm}
  {\Huge\bfseries\color{castblue} Software Requirements Specification\\[0.5em]}
  {\Large swim-times: Automated Swimmer-Times Retrieval System}\\[2em]
  {\large College Area Swim Team (CAST)}\\[0.5em]
  {\normalsize Prepared for: CAST Coaching Staff and Meet Directors}\\[0.5em]
  {\normalsize Prepared by: CAST Technology Working Group}\\[2em]
  \begin{tabular}{ll}
    Document identifier: & CAST-SRS-001                                     \\
    Revision:            & 1.0                                              \\
    Date:                & April 22, 2026                                   \\
    Status:              & Released                                         \\
    Governing standard:  & ISO/IEC/IEEE 29148-2018 (Requirements Engineering) \\
  \end{tabular}
  \vfill
  {\small\textit{This document is prepared in accordance with the principles
  of rational documentation articulated in D.\,L.\,Parnas and P.\,C.\,Clements,
  ``A Rational Design Process: How and Why to Fake It,''
  \textit{IEEE Transactions on Software Engineering}, vol.\,SE-12, no.\,2,
  pp.\,251--257, February 1986.}}
\end{titlepage}

\tableofcontents
\clearpage

% =================================================================
\section{Introduction}
% =================================================================

\subsection{Purpose}

This Software Requirements Specification (SRS) defines the
functional and non-functional requirements for the
\textbf{swim-times} program.  It is the binding agreement between
the requesting stakeholders (CAST coaching staff) and the
implementing developers, and it is the authoritative reference
against which the software shall be verified.

\subsection{Scope}

The software product specified by this document is named
\textbf{swim-times}.  It is a single-binary command-line utility
that:

\begin{itemize}
  \item Accepts swimmer- and event-selection options on the command line.
  \item Resolves each selected swimmer's USA~Swimming \code{PersonKey}
        via the public Sisense JAQL person-search endpoint.
  \item Retrieves all recorded swim times for each selected
        swimmer--event pair from the USA~Swimming Times Elasticube.
  \item Sorts results fastest-first and prints them as either a
        formatted table or comma-separated values.
\end{itemize}

The program does \textbf{not} write to any database, send email,
post to any social channel, or modify any USA~Swimming record.

\subsection{Definitions, acronyms, and abbreviations}

\begin{description}[style=nextline,leftmargin=2.5cm,labelindent=0pt]
  \item[CAST] College Area Swim Team.
  \item[ConOps] Concept of Operations document
                (CAST-CONOPS-001).
  \item[CSV] Comma-Separated Values.
  \item[CWEB] Knuth's literate-programming toolchain for C
              (\code{ctangle} + \code{cweave}).
  \item[ElastiCube] Sisense's columnar in-memory analytics database.
  \item[JAQL] JSON Analytics Query Language.
  \item[LCM] Long-Course Meters.
  \item[PersonKey] USA~Swimming's internal numeric swimmer identifier.
  \item[SCY] Short-Course Yards.
  \item[SDI~AGC] San Diego-Imperial Age-Group Championship.
  \item[SRS] Software Requirements Specification (this document).
\end{description}

\subsection{References}

\begin{enumerate}[label={[\arabic*]}]
  \item ISO/IEC/IEEE 29148-2018, \textit{Systems and software engineering---
        Life cycle processes---Requirements engineering}.
  \item D.\,L.\,Parnas and P.\,C.\,Clements, ``A Rational Design Process:
        How and Why to Fake It,'' \textit{IEEE Trans.\ Software Eng.},
        vol.\,SE-12, no.\,2, pp.\,251--257, Feb.\ 1986.
  \item CAST-CONOPS-001 Rev.\,1.0, \textit{Concept of Operations:
        swim-times}, April 21, 2026.
  \item \code{swim2/swim-times.w} --- CWEB literate program (source of record).
  \item D.\,E.\,Knuth, ``Literate Programming,''
        \textit{The Computer Journal}, vol.\,27, no.\,2,
        pp.\,97--111, 1984.
\end{enumerate}

\subsection{Overview}

Section~2 gives an overall description of the product, its users,
constraints, and assumptions.  Section~3 enumerates the specific
requirements, each tagged with a unique identifier of the form
\code{REQ-X-nn} so that requirements can be traced individually
through verification.  Section~4 specifies how each requirement
shall be verified.

% =================================================================
\section{Overall Description}
% =================================================================

\subsection{Product perspective}

swim-times is a self-contained, end-user command-line tool.
It depends on:

\begin{itemize}
  \item The USA~Swimming Sisense analytics endpoint at
        \texttt{usaswimming.sisense.com} for all data.
  \item \code{libcurl} for HTTPS transport.
  \item A POSIX C runtime for standard library functions.
\end{itemize}

It does not interoperate with any local database, message queue,
or other process.  It is invoked as a child of an interactive
shell or a build script.

\subsection{Product functions}

At a high level, swim-times performs:

\begin{enumerate}
  \item Command-line option parsing (\code{-o}, \code{-e}).
  \item Person-search HTTPS POST to resolve a swimmer's
        \code{PersonKey} from a name substring.
  \item Times-query HTTPS POST per (swimmer, event) pair.
  \item In-memory insertion sort by Sisense \code{SortKey}.
  \item Formatted table or CSV output to standard output.
  \item Diagnostic reporting to standard error.
\end{enumerate}

\subsection{User characteristics}

The program has two user classes:

\begin{description}[style=nextline]
  \item[Coaches and parents.]
        Comfortable with a terminal but not necessarily with
        programming.  They use only the documented command-line
        flags and do not modify the source.
  \item[Maintainers.]
        Familiar with C, CWEB, and \code{libcurl}.  They retangle
        and recompile when the Sisense bearer token rotates or
        when new swimmers, events, or output modes are added.
\end{description}

\subsection{Constraints}

\begin{description}[style=nextline]
  \item[Network.]
        The host running swim-times must have outbound HTTPS access
        to \texttt{usaswimming.sisense.com}.
  \item[Authentication.]
        A valid Sisense bearer token must be embedded in the
        compiled binary.  Token rotation requires a recompile.
  \item[Data source exclusivity.]
        Only the USA~Swimming data hub is consulted.
        Swimcloud is explicitly excluded by stakeholder direction.
  \item[Single-binary distribution.]
        The product is one statically describable executable;
        no installer, no configuration files.
  \item[Literate source of record.]
        The CWEB \code{.w} file is canonical; the generated
        \code{.c} file is a build artefact.
\end{description}

\subsection{Assumptions and dependencies}

\begin{itemize}
  \item USA~Swimming continues to expose the JAQL endpoints
        \code{aPublicIAAaPersonIAAaSearch} and
        \code{aUSAIAAaSwimmingIAAaTimesIAAaElasticube}.
  \item The five output cells per times-query row remain in the
        documented order: time, sort key, meet, date, standard.
  \item No swimmer of interest has more than 200 recorded times in
        any single event (matches the \code{MAX\_TIMES} compile-time
        constant).
\end{itemize}

% =================================================================
\section{Specific Requirements}
% =================================================================

\subsection{Functional requirements}

\subsubsection{Command-line interface}

\begin{description}[style=nextline]
  \item[\req{REQ-F-01}] The program shall accept the option flag
        \code{-o} followed by a comma-separated list of one or more
        of the keywords \code{stella}, \code{kalea}, \code{kenny},
        \code{keith}, \code{fastest}, \code{csv}.
  \item[\req{REQ-F-02}] The program shall accept the option flag
        \code{-e} followed by one or more comma-separated
        USA~Swimming event codes, and shall accept the \code{-e}
        flag repeated multiple times.
  \item[\req{REQ-F-03}] When invoked with no options the program
        shall print a usage message to standard error and exit
        with status~1.
  \item[\req{REQ-F-04}] When invoked with an unrecognised flag the
        program shall print the usage message and exit with status~2.
\end{description}

\subsubsection{Swimmer selection}

\begin{description}[style=nextline]
  \item[\req{REQ-F-10}] When at least one swimmer keyword is
        supplied to \code{-o}, the program shall process only the
        named swimmers.
  \item[\req{REQ-F-11}] When no swimmer keyword is supplied to
        \code{-o} but at least one event is supplied to \code{-e},
        the program shall process every entry in
        \texttt{SWIMMERS[]}: the four Evans/Benavente swimmers plus
        the three age-windowed Ledecky entries
        (\code{ledecky10}, \code{ledecky12}, \code{ledecky14}),
        for seven entries in total.
  \item[\req{REQ-F-12}] The set of selectable swimmers shall be
        exactly: Stella Julianna Evans (\code{stella}),
        Kalea Rose Benavente (\code{kalea}),
        Kenneth Ray Evans (\code{kenny}),
        Keith Santiago Evans (\code{keith}),
        and Katie Genevieve Ledecky (selected via \code{ledecky10},
        \code{ledecky12}, or \code{ledecky14} for the appropriate
        age window).
  \item[\req{REQ-F-13}] When any of the \code{ledecky10},
        \code{ledecky12}, or \code{ledecky14} keywords is
        supplied, the program shall resolve Katie Ledecky's
        \code{PersonKey} via the same person-search path used for
        the other swimmers.
  \item[\req{REQ-F-14}] When the \code{ledecky10} keyword is
        supplied, the program shall restrict output to swims whose
        date falls in the inclusive ISO range
        \code{2006-03-17} through \code{2008-03-16}
        (Katie Ledecky's ninth and tenth years).
  \item[\req{REQ-F-15}] When the \code{ledecky12} keyword is
        supplied, the program shall restrict output to swims whose
        date falls in the inclusive ISO range
        \code{2008-03-17} through \code{2010-03-16}
        (Katie Ledecky's eleventh and twelfth years).
  \item[\req{REQ-F-16}] When the \code{ledecky14} keyword is
        supplied, the program shall restrict output to swims whose
        date falls in the inclusive ISO range
        \code{2010-03-17} through \code{2012-03-16}
        (Katie Ledecky's thirteenth and fourteenth years).
\end{description}

\subsubsection{Event selection}

\begin{description}[style=nextline]
  \item[\req{REQ-F-20}] When one or more event codes are supplied
        to \code{-e}, the program shall query only those events.
  \item[\req{REQ-F-21}] When \code{-e} is omitted, the program
        shall query all thirty-one default events
        (fourteen SCY, seventeen LCM) listed in the swim-times ConOps.
  \item[\req{REQ-F-22}] The program shall support up to
        \code{NUM\_EVENTS} event codes per invocation; additional
        codes beyond that limit shall be silently dropped.
\end{description}

\subsubsection{Data retrieval}

\begin{description}[style=nextline]
  \item[\req{REQ-F-30}] For each selected swimmer the program
        shall resolve a \code{PersonKey} by issuing one HTTPS POST
        to the Sisense person-search JAQL endpoint and matching the
        configured lower-cased substring against the returned
        \code{FullName} field.
  \item[\req{REQ-F-31}] For each (swimmer, event) pair the program
        shall issue one HTTPS POST to the Sisense times-query JAQL
        endpoint, scoped by \code{PersonKey} equals and
        \code{EventCode} equals.
  \item[\req{REQ-F-32}] The program shall request five output
        columns per times-query row in the order: time, sort key,
        meet name, date (\code{YYYYMMDD}), motivational standard.
  \item[\req{REQ-F-33}] The program shall request a maximum of 100
        rows per times query.
\end{description}

\subsubsection{Result processing}

\begin{description}[style=nextline]
  \item[\req{REQ-F-40}] The program shall sort the rows for each
        event by Sisense \code{SortKey} ascending (smallest = fastest)
        before printing.
  \item[\req{REQ-F-41}] When the \code{fastest} keyword is supplied,
        the program shall print only the single fastest row per
        event, suppressing all slower entries.
  \item[\req{REQ-F-42}] The program shall reformat dates from the
        Sisense \code{YYYYMMDD} integer string to \code{YYYY-MM-DD}.
\end{description}

\subsubsection{Output}

\begin{description}[style=nextline]
  \item[\req{REQ-F-50}] In the default (table) mode the program
        shall print, for each event, a heading line, a column
        header, a separator, and one row per recorded time
        containing time, date, motivational standard, and meet name.
  \item[\req{REQ-F-51}] When the \code{csv} keyword is supplied
        the program shall print exactly one CSV header line at the
        start of output and one CSV data line per recorded time.
        Each CSV data line shall contain swimmer name, event code,
        time, date, motivational standard, and meet name in that order.
  \item[\req{REQ-F-52}] When an event has zero recorded times in
        table mode the program shall print the literal
        \code{(no times found)} for that event.
\end{description}

\subsubsection{Error handling}

\begin{description}[style=nextline]
  \item[\req{REQ-F-60}] If a swimmer's \code{PersonKey} cannot be
        resolved the program shall print a diagnostic to standard
        error, free all libcurl global state, and exit with
        status~1.
  \item[\req{REQ-F-61}] If an individual times query fails, the
        program shall print a diagnostic to standard error and
        continue with the next event without aborting.
\end{description}

\subsubsection{Database persistence}

\begin{description}[style=nextline]
  \item[\req{REQ-F-70}] When the \code{store} keyword is supplied,
        the program shall write every emitted row to the local
        Postgres database \texttt{swim-times} programmatically via
        the \texttt{libpq} client library, using a prepared
        \texttt{INSERT \dots\ ON~CONFLICT~DO~NOTHING} per row.
        \code{store} composes with \code{csv}: stdout still
        receives the CSV stream while the rows are persisted.
  \item[\req{REQ-F-71}] When the \code{offline} keyword is
        supplied, the program shall read swim rows from the local
        Postgres \texttt{swim-times} database and emit them in the
        same table or CSV form as the online path, without
        touching the network.
  \item[\req{REQ-F-72}] Under \code{offline}, the program shall
        not initialise libcurl and shall not issue any HTTP
        request.
  \item[\req{REQ-F-73}] The \texttt{swim\_times} schema shall
        enforce a UNIQUE constraint on
        \code{(swimmer, event, swim\_time, swim\_date, meet)} so
        that re-running a \code{store} invocation does not
        introduce duplicate rows.
  \item[\req{REQ-F-74}] If the libpq connection cannot be
        established for \code{store}, the program shall print a
        single warning to standard error and continue with the
        run, without persisting any rows.  Under \code{offline} a
        connection failure is fatal and exits with status~1.
  \item[\req{REQ-F-75}] When \code{store} is rerun against a
        \texttt{swim\_times} table that already contains the rows
        being inserted, the program shall recover gracefully: each
        duplicate row shall be silently absorbed by the
        \texttt{ON~CONFLICT~DO~NOTHING} clause without aborting
        the transaction, the process shall exit~0, and no
        diagnostic shall be written to standard error.
        (This realises requirement~\#4 of \texttt{requirements4.txt}.)
\end{description}

\subsection{External interface requirements}

\subsubsection{User interfaces}

\begin{description}[style=nextline]
  \item[\req{REQ-I-01}] All user interaction shall be via
        command-line arguments and process exit status; the
        program shall not read from standard input.
\end{description}

\subsubsection{Software interfaces}

\begin{description}[style=nextline]
  \item[\req{REQ-I-10}] The program shall use the
        \code{libcurl} ``easy'' interface for all HTTP transport.
  \item[\req{REQ-I-11}] The program shall send the headers
        \code{Content-Type: application/json} and
        \code{Authorization: Bearer \{token\}} on every JAQL
        request.
  \item[\req{REQ-I-12}] The program shall reach Postgres
        \emph{programmatically} through the \texttt{libpq} client
        library --- never by spawning \texttt{psql} or any other
        external script (this realises requirement~\#3 of
        \texttt{requirements4.txt}).  The connection string shall
        be taken from the \texttt{SWIM\_TIMES\_PGCONNINFO}
        environment variable when set, otherwise the default
        keyword string \texttt{dbname=swim-times user=postgres
        host=localhost port=5432}.
\end{description}

\subsubsection{Communication interfaces}

\begin{description}[style=nextline]
  \item[\req{REQ-I-20}] All network communication shall use HTTPS
        (TLS) to the host \texttt{usaswimming.sisense.com}.
\end{description}

\subsection{Performance requirements}

\begin{description}[style=nextline]
  \item[\req{REQ-P-01}] A full refresh of all seven
        \texttt{SWIMMERS[]} entries across all thirty-one default
        events ($7 \times (1 + 31) = 224$ HTTP requests, including
        the three Ledecky age-window entries that share an
        underlying person but each issue their own person-search
        and event queries) shall complete within five minutes on a
        residential broadband connection.
  \item[\req{REQ-P-02}] A single-swimmer single-event query
        (\code{-o stella -e "100 FR SCY"}) shall complete within
        ten seconds under normal Sisense response time.
  \item[\req{REQ-P-03}] An offline single-swimmer single-event
        query (\code{-o stella,offline -e "100 FR SCY"}) shall
        complete within two seconds against a local Postgres
        instance carrying typical row counts.
\end{description}

\subsection{Design constraints}

\begin{description}[style=nextline]
  \item[\req{REQ-D-01}] The software shall be implemented in C
        conformant to ISO/IEC 9899:1999 (C99) or later.
  \item[\req{REQ-D-02}] The canonical source shall be a CWEB
        literate program (\code{swim-times.w}); the
        \code{.c} file shall be a build artefact only.
  \item[\req{REQ-D-03}] No source-code module (named CWEB chunk)
        shall exceed twenty-four lines, in keeping with the
        existing literate-programming style of the codebase.
  \item[\req{REQ-D-04}] No third-party dependency other than
        \code{libcurl}, \code{libpq} (the Postgres C client
        library, added by requirements4.txt~\#3), and the standard
        C library shall be linked into the binary.  Transitive
        TLS, compression, and authentication libraries pulled in
        by \code{libcurl} or \code{libpq} are also permitted.
  \item[\req{REQ-D-10}] The Postgres schema for persisted swim
        rows shall be defined in \texttt{schema.sql} adjacent to
        the source.  The table \texttt{swim\_times} shall expose
        the seven application columns listed in REQ-F-32 plus an
        identity primary key and a \texttt{fetched\_at} timestamp,
        and shall declare a UNIQUE constraint per REQ-F-73.
\end{description}

\subsection{Software system attributes}

\begin{description}[style=nextline]
  \item[\req{REQ-A-01} (Reliability)] The program shall produce
        identical output for identical Sisense responses on
        repeated invocations.
  \item[\req{REQ-A-02} (Maintainability)] All Sisense endpoint
        URLs and bearer tokens shall be declared as named
        constants in clearly identifiable CWEB chunks.
  \item[\req{REQ-A-03} (Portability)] The program shall build
        and run unmodified on any POSIX.1-2008 system with
        \code{libcurl} 7.0 or later installed.
  \item[\req{REQ-A-04} (Self-documentation)] The CWEB source
        shall, when processed by \code{cweave} and \code{pdftex},
        produce a human-readable typeset document describing both
        the why and the how of the implementation.
\end{description}

% =================================================================
\section{Verification}
% =================================================================

Each requirement shall be verified by one of the following methods:

\begin{description}[style=nextline]
  \item[Inspection (I).]
        Visual examination of source, output, or documentation.
  \item[Demonstration (D).]
        Execution of the program with selected inputs and
        observation of the resulting output.
  \item[Test (T).]
        Execution of an automated test that compares actual
        output to an expected reference.
  \item[Analysis (A).]
        Reasoning about the program (e.g., line-count audit
        of CWEB chunks for \code{REQ-D-03}).
\end{description}

\medskip
\noindent
The verification matrix below assigns a primary method to each
requirement.  Test cases and demonstration scripts shall be
documented separately in the swim-times Verification Plan.

\begin{center}
\begin{longtable}{ll}
  \toprule
  \textbf{Requirement} & \textbf{Method} \\
  \midrule\endhead
  REQ-F-01 \ldots REQ-F-04 & Demonstration \\
  REQ-F-10 \ldots REQ-F-16 & Demonstration \\
  REQ-F-20 \ldots REQ-F-22 & Demonstration \\
  REQ-F-30 \ldots REQ-F-33 & Inspection + Demonstration \\
  REQ-F-40 \ldots REQ-F-42 & Test \\
  REQ-F-50 \ldots REQ-F-52 & Demonstration \\
  REQ-F-60 \ldots REQ-F-61 & Demonstration \\
  REQ-F-70 \ldots REQ-F-75 & Test (DB-gated) \\
  REQ-I-01 \ldots REQ-I-20 & Inspection \\
  REQ-P-01, REQ-P-02, REQ-P-03 & Demonstration (timed) \\
  REQ-D-01, REQ-D-02, REQ-D-04 & Inspection \\
  REQ-D-03                 & Analysis (CWEB chunk line-count audit) \\
  REQ-D-10                 & Inspection of \texttt{schema.sql} \\
  REQ-A-01                 & Test \\
  REQ-A-02, REQ-A-04       & Inspection \\
  REQ-A-03                 & Demonstration (build on second POSIX host) \\
  \bottomrule
\end{longtable}
\end{center}

\vfill
\begin{center}
  \small\textit{End of document CAST-SRS-001 Rev.\,1.0}
\end{center}

\end{document}
